% ============================================================
% Plantilla: Artículo con normas APA 7.ª edición (formato de estudiante)
% Las citas y referencias las genera biblatex-apa (Recompilar ejecuta biber automáticamente).
%   \textcite{clave}  ->  Autor (año)
%   \parencite{clave} ->  (Autor, año)
% ============================================================
\documentclass[12pt,a4paper]{article}

\usepackage[utf8]{inputenc}
\usepackage[T1]{fontenc}
\usepackage{mathptmx}                 % Times New Roman, como pide APA
\usepackage[spanish,es-tabla]{babel}
\usepackage[margin=2.54cm,headheight=15pt]{geometry}  % márgenes de 1 pulgada
\usepackage{setspace}
\usepackage{csquotes}
\usepackage{booktabs}
\usepackage{caption}
\usepackage[style=apa,backend=biber]{biblatex}
\usepackage[hidelinks]{hyperref}

\addbibresource{referencias.bib}
\DeclareLanguageMapping{spanish}{spanish-apa}

\doublespacing
\setlength{\parindent}{1.27cm}

% Número de página arriba a la derecha
\usepackage{fancyhdr}
\pagestyle{fancy}
\fancyhf{}
\fancyhead[R]{\thepage}
\renewcommand{\headrulewidth}{0pt}

% Títulos con los niveles de APA (centrado en negrita, izquierda en negrita, izquierda en negrita cursiva)
\makeatletter
\renewcommand{\section}{\@startsection{section}{1}{0pt}{1em}{0.5em}{\centering\normalfont\bfseries}}
\renewcommand{\subsection}{\@startsection{subsection}{2}{0pt}{1em}{0.5em}{\normalfont\bfseries}}
\renewcommand{\subsubsection}{\@startsection{subsubsection}{3}{0pt}{1em}{0.5em}{\normalfont\bfseries\itshape}}
\makeatother
\setcounter{secnumdepth}{0}

\captionsetup{labelfont=bf, labelsep=newline, singlelinecheck=false, textfont=it}

\begin{document}

% ---------- Portada de estudiante (APA 7) ----------
\begin{titlepage}
  \centering
  \vspace*{3cm}
  {\bfseries Título del Trabajo en Mayúsculas y Minúsculas\par}
  \vspace{2em}
  Nombre del Estudiante\par
  Departamento o Facultad, Nombre de la Universidad\par
  Código y Nombre de la Asignatura\par
  Nombre del Docente\par
  \today\par
\end{titlepage}

\setcounter{page}{2}

% ---------- Cuerpo ----------
\begin{center}
  \textbf{Título del Trabajo en Mayúsculas y Minúsculas}
\end{center}

El texto comienza aquí con sangría en la primera línea de cada párrafo. Según
\textcite{apa2020}, el estilo APA busca claridad y precisión en la comunicación
científica. Las citas parentéticas se escriben así \parencite{lamport1994}; con dos
autores, dentro del paréntesis se unen con «\&» \parencite{goossens2004}.

\section{Método}
Describe a los participantes, los instrumentos y el procedimiento.

\subsection{Participantes}
Indica el número de participantes y sus características.

\subsubsection{Criterios de inclusión}
Detalla los criterios utilizados.

\section{Resultados}
La Tabla~\ref{tab:descriptivos} presenta los estadísticos descriptivos.

\begin{table}[htbp]
  \caption{Estadísticos descriptivos de las variables}
  \label{tab:descriptivos}
  \begin{tabular}{lcc}
    \toprule
    Variable & \textit{M} & \textit{DE} \\
    \midrule
    Edad & 21.4 & 2.3 \\
    Promedio académico & 8.7 & 0.9 \\
    \bottomrule
  \end{tabular}
\end{table}

\section{Discusión}
Interpreta los resultados a la luz de la literatura \parencite{knuth1984}.

\newpage
\printbibliography[title=Referencias]

\end{document}
